\documentclass[11pt]{article}


\setlength{\parindent}{0pt}
\usepackage{xltxtra}
\usepackage{hyperref}
\hypersetup{hidelinks}
\usepackage{url}
\urlstyle{tt}
\usepackage{xcolor}
\definecolor{CVBlue}{RGB}{23,110,191}
\usepackage{calc}
\usepackage{graphicx}
\usepackage{tikz}
\usetikzlibrary{calc}
\usepackage{fontspec}
\usepackage{xeCJK}
\usepackage{enumitem}
\CJKsetecglue{} %% 取消中文与数字之间的间隙


%% 主文档字体设置
\setmainfont[
    Path = fonts/Main/,
    Extension = .otf,
    BoldFont = texgyretermes-bold.otf, % 加粗字体
]{texgyretermes-regular.otf} % 正文字体

% 中文字体设置
\setCJKmainfont[
    Path = fonts/hansans/,
    Extension = .ttf,
    BoldFont = NotoSansSC-Bold.ttf, % 加粗字体
]{NotoSansSC-Regular.otf} % 正文字体


\usepackage{relsize}
\usepackage{xspace}

% 使用 fontawesome（部分图标）
\usepackage{fontawesome} 

% A4纸，上下左右边距
\usepackage[
    a4paper,
    left=1.2cm,
    right=1.2cm,
    top=1.5cm,
    bottom=1cm,
    nohead
]{geometry}

\renewcommand{\baselinestretch}{1.5} % 行间距设为1.5

\usepackage{titlesec}
\usepackage{enumitem}
\setlist{noitemsep} % 取消列表项间的额外间距
%\setlist{nosep} % 取消所有垂直间距
\setlist[itemize]{topsep=0.25em, leftmargin=*}
\setlist[enumerate]{topsep=0.25em, leftmargin=*}

% --- 用于控制【不同项目之间】的垂直距离 ---
\newlength{\interProjectSpacing}
\setlength{\interProjectSpacing}{0.9em} % <--- 在此调整项目之间的距离
\newcommand{\projectsep}{\vspace{\interProjectSpacing}}

% --- 用于控制【项目标题】与下方【项目描述】的距离 ---
\newlength{\intraProjectTitleSep}
\setlength{\intraProjectTitleSep}{0.4em} % <--- 在此调整标题和描述的距离
\newcommand{\titlebreak}{\\[\intraProjectTitleSep]}

% --- 用于控制【项目描述】与下方【要点列表】的距离 ---
\newlength{\intraProjectListTopSep}
\setlength{\intraProjectListTopSep}{0.2em} % <--- 在此调整描述和列表的距离

% =======================================================================


\titleformat{\section}         % 定制 \section 命令 
{\large\bfseries\raggedright} % 将 section 标题设置为大号、粗体且左对齐
{}{0em}                      % 可用于为所有 section 添加前缀（如“章节...”）
{}                           % 可用于在标题前插入代码
[{\color{CVBlue}\titlerule}]  % 在标题后插入一条横线
\titlespacing*{\section}{0cm}{*1.6}{*1.2}



\begin{document}
\pagenumbering{gobble}

  %%%% 利用tikz来定位学校Logo，这里只在第一页显示
  \begin{tikzpicture}[remember picture, overlay] 
    \node[anchor = north west] at ($(current page.north west)+(0.5cm,+1.0cm)$) {\includegraphics[height=6cm]{zju.png}};
  \end{tikzpicture}%
\centerline{\LARGE\bfseries{金诣涵}} 

\centerline{\normalsize{\faPhone\ 15958976766 \quad \faEnvelopeO\ \href{mailto:3230103210@zju.edu.cn}{3230103210@zju.edu.cn}}} 
    
\section{\makebox[\widthof{\faGraduationCap}][c]{\color{CVBlue}\faGraduationCap}\ 教育背景}    
\textbf{浙江大学} \hfill 2023.9 -- 至今\\[0.5em] % 标题和正文间加一点距离
电子科学与技术\quad 大三 
\begin{itemize}[nosep]
    \item 总平均绩点：4.52 主修平均绩点：4.48
    \item 总学分：151.5
    \item 成绩排名： 18/86
    \item 外语成绩： CET6（612），IELTS （7.0） 
\end{itemize}

\section{\makebox[\widthof{\faUsers}][c]{\color{CVBlue}\faUsers}\ 科研竞赛经历}

% --- 第一个项目 ---
% 将标题行末尾的 \\ 替换为 \titlebreak 命令
\textbf{全国大学生电子设计竞赛G题：电路模型探究装置} \hfill 2025.08 \titlebreak
项目描述：设计并制作题目所给指标的RC有源滤波电路和满足以下要求的电路模型探究装置：（1）可产生指定频率幅度的正弦信号，（2）可学习测评现场提供由RLC元件组成的“未知模型电路”的类型（3）将信号发生器接入探究装置和“未知模型电路”的输入端口，能根据信号发生器的输出信号推理生成与“未知模型电路”相同的输出信号。要求探究装置和“未知模型电路”的输出信号相比波形无失真、在示波器上能连续同频稳定显示（相位无要求），两者峰峰值相对误差绝对值不大于10%。 
% 在 itemize 的选项中，使用 topsep=\intraProjectListTopSep 来控制上边距
\begin{itemize}[nosep, topsep=\intraProjectListTopSep]
    \item \textbf{已知模型电路}：运用NE5532搭建二阶SK低通滤波器，并搭建由OPA828高速运放对滤波后的信号放大的放大电路。
    \item \textbf{DDS输出放大}：考虑到单片机驱动DDS输出的正弦波幅度较小且最高频率高于1MHz，选用高速低噪放大器OPA847芯片搭建放大电路。
    \item \textbf{锁相电路}：为了使探究装置与未知模型电路输出波形稳定同频显示且波形之间无漂移，需要使得探究装置的输出信号与信号发生器的时钟信号相关，采用CD4046搭建锁相环并使用74HC161N搭建分频器构成闭环的锁相倍频电路，实现将信号发生器信号转化为16倍频时钟信号。
\end{itemize}

% 使用 \projectsep 命令来分隔两个项目
\projectsep

% --- 第二个项目 ---
\textbf{基于深度学习和最大似然估计的波达方向估计算法研究} \hfill 2025.03 -- 2026.08 \titlebreak
项目描述：现有的高分辨率DOA估计算法主要基于数学模型，在非理想环境下性能会严重下降，而现有的传统深度学习方法在高信噪比情况下存在固有的精度瓶颈。该项目通过使用生成式模型学习数据的概率分布，来达到兼顾深度学习的强鲁棒性与MLE的渐近统计最优性的目的。
\begin{itemize}[nosep, topsep=\intraProjectListTopSep]
    \item \textbf{前期（2025.03 -- 2025.09）}：配置搭建了所需的Pytorch环境，学习并确定整体研究思路，根据研究所需场景构建含阵列误差，非高斯噪声的一维均匀线阵仿真数据集。
    \item \textbf{中期（2025.10 -- 2026.03）}：以归一化流Normalizing Flow为基础架构，在此基础上搭建学习MLE的生成式模型，提出生成式最大似然（GMLE）算法，并通过仿真对比验证GMLE算法与Root-MUSIC及CNN在存在固定阵列误差以及非高斯干扰噪声等复杂场景下的RMSE表现。
    \item \textbf{后期（2026.04 -- 2026.05）}：对整体算法思路，源代码以及数据集进行整理，撰写，打磨并产出了一篇会议论文，现已被CIE Radar 2026国际会议录用。
\end{itemize}

\section{\makebox[\widthof{\faCogs}][c]{\color{CVBlue}\faCogs}\ 部分课程成绩}
\begin{itemize}[nosep]
    \item \textbf{集成电路原理与设计：} 95/100
    \item \textbf{数字信号处理：} 94/100
    \item \textbf{射频电路与系统：} 98/100
    \item \textbf{信息电子学物理基础：} 90/100
\end{itemize}
\section{\makebox[\widthof{\faGraduationCap}][c]{\color{CVBlue}\faList}\ 荣誉情况}
\begin{itemize}
    \item 校设奖学金 浙江大学三等奖学金 \hfill 2023-2024学年
    \item 校设奖学金 浙江大学三等奖学金 \hfill 2024-2025学年
    \item 校设奖学金 国家级人才培养基地奖学金国家级人才培养基地三等奖学金 \hfill 2024-2025学年
\end{itemize}
\section{\makebox[\widthof{\faGraduationCap}][c]{\color{CVBlue}\faList}\ 获奖情况}
\begin{itemize}
    \item 浙江大学“TP-LINK 杯”第 30 届大学生电子设计竞赛校级三等奖 \hfill 2024年
    \item 浙江大学“TP-LINK 杯”第 31 届大学生电子设计竞赛校级一等奖 \hfill 2025年
    \item “TI 杯”全国大学生电子设计竞赛（浙江赛区）省级一等奖及推荐国家奖 \hfill 2025年
    \item “TI 杯”全国大学生电子设计竞赛（国奖测评）全国一等奖 \hfill 2025年
\end{itemize}
\end{document}
